\section{反垄断法}
\begin{enumerate}
  \item 分值可能突破，2 题。
  \item 范围。
        \begin{enumerate}
          \item \yi 境内。\er 在境外，但对境内有影响。\al{微软收购暴雪}。
          \item 不适用。知识产权行为（天然排他）、农业。
        \end{enumerate}
  \item 行政主体在制定涉及市场经济规定时，应当进行\bs{公平竞争审查}。
  \item 反垄断机构。
        \begin{enumerate}
          \item 反垄断委员会 \yjt 组织、协调、指导。
          \item 反垄断\zd{执法}机构。国家市场监督管理总局/\zd{市监局 \yjt 授权省}。
                \begin{enumerate}
                  \item 商务部。主管经营者集中。
                  \item 发改委。主管价格垄断。
                  \item 工商总局。主管其他垄断行为。
                \end{enumerate}
        \end{enumerate}
  \item 3 种垄断行为。\bs{垄断协议、滥用市场支配地位、经营者集中}。
        \begin{enumerate}
          \item 垄断协议。\zd{横向、纵向、轴辐}。
                \begin{enumerate}
                  \item 横向。具有\zd{竞争关系}的主体。竞争关系 = 类似行为。
                        \begin{enumerate}
                          \item 变更商品\zd{价格}。\al{联手涨价}。
                          \item 限制生产 \hz 销售\zd{数量}。
                          \item 分割销售 \hz 原材料采购\zd{市场}。\al{中石油、中石化划江而治}。
                          \item \zd{限制开发}/购买新技术、新材料。\al{环保技术不推广默契}。
                          \item 联合抵制。
                        \end{enumerate}
                  \item 纵向。经营者与交易相对人，\zd{上下游关系}。更具隐蔽性。
                        \begin{enumerate}
                          \item 固定转售价格。\al{全国统一零售价 \cuo。建议零售价 \dui}。
                          \item 限定最低价格。\al{每个轮胎的价格不得低于 2000 元}。
                          \item \zd{安全港制度}。市场份额低于标准。
                        \end{enumerate}
                  \item \bs{轴辐}。大概率会考到。
                        \begin{enumerate}
                          \item \zd{组织型}。经营者不得\zd{组织其他经营者达成}垄断协议。
                          \item \zd{帮助型}。经营者不得为其他经营者达成垄断协议\zd{提供实质性的帮助}。
                        \end{enumerate}
                  \item 例外。\zx{不排除、限制竞争}（\zd{没有坏效果}）。\al{疫情期间口罩卖 2 毛}。
                  \item 豁免。目的正当，效果积极（不是出于私心）。常识判断即可。\al{一致对外，保障外贸}。
                  \item 2 个新增。\yi \zd{加倍处罚}。特别严重、恶劣、严重后果。\er \zd{对经营负责人的处罚}。
                  \item \bs{宽恕/宽大制度只针对垄断协议}。
                \end{enumerate}
          \item 滥用市场支配地位。
                \begin{enumerate}
                  \item \bs{有市场支配地位 \bdy 违法}。\al{甲公司具有市场支配地位，应介入调查并予以处罚 \cuo}。
                  \item 认定因素。
                        \begin{enumerate}
                          \item 市场\zd{份额}巨大（\bs{70\%}），具有\zd{定价权}。
                          \item \zd{控制}销售市场 \hz 原材料采购市场的能力。
                          \item 经营者\zd{财力和技术}条件、交易市场\zd{依赖}程度、\zd{进入}市场难易程度。
                        \end{enumerate}
                  \item \bs{推定规则}。1个 1/2、2 个 2/3、3 个 3/4。\zd{不到 1/10 \yjt 小跟班}。\al{甲 1/2，乙 1/6，丙 1/12}。
                  \item 行为。
                        \begin{enumerate}
                          \item 垄断价格。\al{知网}。
                          \item 强制交易。\al{阿里、美团二选一}。
                          \item 低价倾销、拒绝交易、搭售行为、差别待遇。
                        \end{enumerate}
                  \item 数字经济领域反垄断。\al{二选一、大数据杀熟}。
                  \item 法律责任。
                        \begin{enumerate}
                          \item 停止违法行为、没收违法所得。
                          \item \bs{上年度销售额（不是利润额）1\% - 10\% 的罚款}。
                          \item \zd{新增加倍处罚（2 - 5 倍）}。特别严重、恶劣、严重后果。\al{追究刑事责任 \cuo}。
                        \end{enumerate}
                \end{enumerate}
          \item 经营者集中。取得控制权。复议前置但不终局。
                \begin{enumerate}
                  \item \zd{交易规模}达到国务院规定的申报标准。\bdy 公司法的合并。\al{可口可乐收购汇源果汁}。
                  \item \zx{事先申报 \hz 主动调查}。\al{事后申报 \cuo}。例外。\zx{母子、兄弟公司整合}。\al{甲控股乙、丙，乙收购丙股份}。
                  \item 分级分类审查。
                  \item 初审 \yjt 进一步审查。30 + 90 + 60。
                  \item \zd{停表制度}（中止计算期限）。少交材料；新情况、事实；经营者提出。
                  \item 股份恢复到集中前的状态。\al{解除独家版权}。
                \end{enumerate}
          \item \sjt \zx{存在民事责任}。\al{可对其提起民事诉讼，要求赔偿损失}。
        \end{enumerate}
  \item 行政垄断。\zd{强制交易、不平等待遇、准入}。\zx{不能直接处罚} \yjt 没有行政隶属关系。
  \item 措施。
        \begin{enumerate}
          \item 承诺 \yjt 暂停调查。做到了 \yjt 结束。
          \item \zd{约谈。新增}。
          \item \zd{设区的市以上检察院}进行反垄断公益诉讼。
          \item 行政处罚\zd{纳入公信}，向社会公布。
        \end{enumerate}
\end{enumerate}
